
%%%%%%%%%%%%%%%%%%%%%%%%%%%%%%%%%%%%%%%%%%%%%%%%%%%%%%%%%%%
%% Inicio del capítulo de fundamentos
%%%%%%%%%%%%%%%%%%%%%%%%%%%%%%%%%%%%%%%%%%%%%%%%%%%%%%%%%%%

\chapter{Fundamentos}

En este capítulo se introducen los conceptos teóricos y tecnológicos necesarios para comprender el resto del trabajo: los modelos de lenguaje de gran escala (LLMs) y su funcionamiento, el estándar BPMN y su notación básica, y la clasificación PERVAL para el análisis del valor percibido por el cliente.

%%---------------------------------------------------------
\section{Modelos de lenguaje de gran escala (LLMs)}
%%---------------------------------------------------------

Los modelos de lenguaje de gran escala (\textit{Large Language Models,} LLMs)  constituyen una categoría de modelos fundacionales (\textit{foundation models) basados en}  aprendizaje profundo y entrenados sobre grandes volúmenes de datos textuales. Su  objetivo  es aprender  representaciones estadísticas del lenguaje  que  permitan predecir y generar  secuencias de texto coherentes, así como realizar diversas tareas de procesamiento del lenguaje natural a partir de instrucciones expresadas en lenguaje natural  \cite{bommasani2021}. 

La mayoría de los LLMs actuales se basan en la arquitectura \textit{transformer}, propuesta por Vaswani et al. \cite{vaswani2017}, que sustituye las redes  neuronales recurrentes empleadas tradicionalmente en el procesamiento del lenguaje natural por un mecanismo de \textit{self-attention} (auto-atención). Este mecanismo permite al modelo ponderar, para cada palabra o fragmento de texto (\textit{token}) de una secuencia, la relevancia del resto de los elementos de dicha secuencia,  facilitando así la capturar  eficiente y paralelizable de dependencias de largo alcance entre palabras.

El entrenamiento de estos modelos se organiza habitualmente en dos fases. En primer lugar, una fase de \textit{pre-training} (preentrenamiento), en la que el modelo aprende a predecir el siguiente \textit{token} de una secuencia a partir de grandes corpus de texto no etiquetado. Mediante este proceso,  el modelo desarrolla representaciones internas que codifican regularidades estadísticas del lenguaje, hechos y patrones de razonamiento presentes en los datos de entrenamiento. En segundo lugar, opcionalmente,  puede llevarse a cabo una fase de \textit{fine-tuning} (ajuste fino), en la que el modelo se especializa en tareas concretas mediante conjuntos de datos más reducidos y específicos.

Diversos estudios han mostrado que, a medida que aumenta la escala de estos modelos, emergen capacidades que no se manifiestan de forma evidente en modelos más pequeños, tales como la resolución de nuevas tareas a partir de unos pocos ejemplos (\textit{few-shot learning), la adaptación a dominios específicos y la capacidad de realizar razonamientos complejos utilizando únicamente la información proporcionada en el contexto de entrada (in-context learning) \cite{brown2020, wei2022}. En particular,} Brown \textit{et al}. \cite{brown2020} demostraron que modelos de gran escala pueden  resolver  tareas  para las que no han sido entrenados explícitamente, utilizando únicamente ejemplos e instrucciones incluidos en el propio \textit{prompt}.

La forma habitual de interactuar con un LLM consiste en proporcionarle una instrucción en lenguaje natural, denominada \textit{prompt}, que describe la tarea a realizar y, opcionalmente, el formato esperado de la respuesta. La calidad y precisión de la salida generada depende en gran medida de cómo se formula dicho \textit{prompt}, así como del contexto y de  las instrucciones adicionales  proporcionadas al modelo, dando lugar a la disciplina conocida como \textit{prompt engineering}.

A pesar de sus capacidades, los LLMs también presentan limitaciones importantes. Entre ellas destacan la generación de información incorrecta o no fundamentada (\textit{hallucinations), la sensibilidad a pequeñas variaciones en la formulación de las instrucciones y la dificultad para garantizar la corrección semántica de los artefactos generados \cite{ji2023, feuerriegel2024}. Estas limitaciones resultan especialmente relevantes cuando los modelos se utilizan en tareas de ingeniería del software y modelado de procesos de negocio, en las que la precisión y la consistencia de los artefactos producidos constituyen requisitos fundamentales.}

Estas capacidades   y limitaciones son especialmente relevantes para el presente TFM. Por una parte,  los LLMs ofrecen la posibilidad de interpretar descripciones de procesos de negocio expresadas en lenguaje natural y genere, a partir de ellas, representaciones estructuradas, como modelos BPMN. Por otra parte, su naturaleza conversacional permite mantener un diálogo iterativo con el usuario para refinar progresivamente dichos modelos a partir de nuevas instrucciones  y restricciones. La aplicación de estas capacidades al ámbito de la gestión de procesos de negocio ha sido objeto de  creciente atención en investigaciones  recientes \cite{kourani2025evaluating, kourani2024process, klievtsova2023, abb2023, bise2025bpmn, fillies2024}, cuyo  principales avances y limitaciones se analizan con mayor  detalle en el capítulo dedicado al Estado del Arte.

%%---------------------------------------------------------
\section{BPMN: notación y modelado de procesos de negocio}
%%---------------------------------------------------------

La Gestión de Procesos de Negocio (\textit{Business Process Management, BPM) se ha consolidado como una disciplina fundamental para las organizaciones que buscan mejorar su eficiencia operativa, incrementar la calidad de sus servicios y apoyar iniciativas de transformación digital. En este contexto, la modelización de procesos constituye una actividad esencial para comprender, documentar, analizar y rediseñar la forma en que las organizaciones crean y entregan valor \cite{dumas2018fundamentals, weske2019}.}

BPMN  proporciona una notación común  y estandarizada que facilita la comunicación entre analistas de negocio, expertos del dominio, desarrolladores de sistemas y demás partes interesadas involucradas en el diseño y ejecución de los  procesos  \cite{omg2011bpmn, dumas2018fundamentals}. Su objetivo principal es ofrecer una notación gráfica fácilmente entendible por todos los usuarios  con distintos perfiles técnicos y de negocio , al tiempo que proporciona un nivel de expresividad suficiente para modelar procesos de diferente complejidad \cite{dumas2018fundamentals}.

A pesar de su amplia adopción en entornos académicos e industriales, la construcción manual de modelos BPMN continúa siendo una actividad cognitivamente exigente. La elaboración de modelos de calidad requiere conocimientos especializados sobre la notación, comprensión del dominio del problema y experiencia en modelado de procesos. Asimismo, la complejidad de los modelos puede aumentar considerablemente cuando intervienen múltiples participantes, excepciones, rutas alternativas de ejecución o elevados niveles de detalle \cite{dumas2018fundamentals, mendling2010}.

Un diagrama BPMN (\textit{Business Process Diagram}, BPD) se construye a partir de cuatro categorías básicas de elementos:

\begin{itemize}
    \item[$\bullet$] \textbf{Objetos de flujo} (\textit{flow objects}): constituyen los elementos principales del proceso e incluyen los \textit{events} (eventos), representados mediante círculos y  utilizados para  indicar algo que sucede durante la ejecución del proceso (eventos de inicio,  intermedios o de fin); las \textit{activities} (actividades), representadas mediante rectángulos de esquinas redondeadas y  utilizadas para  representar el trabajo realizado; y los \textit{gateways} (compuertas), representados mediante rombos y  encargadas de  controlar la divergencia y convergencia del flujo  de ejecución, por ejemplo mediante compuertas exclusivas, paralelas o rutas inclusivas.
    
    \item[$\bullet$] \textbf{Objetos de conexión} (\textit{connecting objects}): incluyen los \textit{sequence flows} (flujos de secuencia), que indican el orden de ejecución de las actividades; los \textit{message flows} (flujos de mensaje), que representan el intercambio de información entre participantes; y las \textit{associations} (asociaciones), que vinculan artefactos a otros elementos del diagrama.
    
    \item[$\bullet$] \textbf{Carriles} (\textit{swimlanes}):  comprenden los \textit{pools} y \textit{lanes}, que permiten organizar gráficamente las actividades  de acuerdo con  el participante, la organización o el rol responsable de su ejecución.
    
    \item[$\bullet$] \textbf{Artefactos} (\textit{artifacts}):  incluyen elementos adicionales, como los \textit{data objects} (objetos de datos), los \textit{groups} (grupos) y  las \textit{annotations} (anotaciones), que aportan información complementaria sin afectar  el comportamiento del  flujo de proceso.
\end{itemize}

Además de su representación gráfica, la especificación BPMN 2.0 define un formato de intercambio basado en XML que permite representar, almacenar e intercambiar modelos de proceso entre distintas herramientas de modelado \cite{omg2011bpmn}. Esta  característica es fundamental para el presente TFM, ya que el LLM integrado en la extensión genera y modifica los modelos de proceso directamente en este formato XML, que posteriormente es interpretado y renderizado por bpmn.io para su visualización y edición dentro del entorno de modelado.

Esta característica resulta particularmente relevante para el presente TFM. La solución desarrollada utiliza una instancia local de bpmn.io como entorno de modelado y aprovecha la representación XML de BPMN como mecanismo de intercambio entre el LLM y el editor gráfico. En particular, el modelo de lenguaje genera propuestas de modelos de proceso expresadas en XML BPMN, las cuales son posteriormente interpretadas y renderizadas por bpmn.io para su visualización, edición y refinamiento iterativo dentro del entorno de modelado.

%%---------------------------------------------------------
\section{La clasificación PERVAL}
%%---------------------------------------------------------

El concepto de valor percibido ha sido ampliamente estudiado en las áreas de marketing, comportamiento del consumidor y gestión de servicios, debido a su importancia para comprender cómo los clientes evalúan productos, servicios y experiencias de consumo. En términos generales, el valor percibido puede entenderse como la valoración global que realiza un cliente sobre la utilidad de un producto o servicio, considerando el equilibrio entre los beneficios obtenidos y los sacrificios realizados para acceder a él, tales como costes económicos, esfuerzo, tiempo o riesgos percibidos \cite{zeithaml1988}.

Entre las distintas aproximaciones propuestas en la literatura, la clasificación PERVAL (\textit{Perceived Value}) fue propuesta por Sweeney y Soutar \cite{sweeney2001perval} como una escala multidimensional  destinada a medir el valor percibido por el cliente en relación con un producto y servicio. A diferencia de enfoques unidimensionales centrados exclusivamente en aspectos económicos o funcionales, PERVAL asume que la percepción de valor constituye un fenómeno complejo y multidimensional, en el que intervienen simultáneamente componentes funcionales, emocionales, sociales y económicos. Esta perspectiva ha favorecido su amplia utilización en investigaciones relacionadas con el comportamiento del consumidor, la evaluación de servicios y el análisis de experiencias de usuario.

La escala PERVAL identifica cuatro dimensiones principales del valor percibido:

\begin{itemize}
    \item[$\bullet$] \textbf{Valor emocional} (\textit{emotional value}): utilidad derivada de los sentimientos o estados afectivos que  el producto o servicio genera en el cliente, tales como satisfacción, tranquilidad, diversión o confianza.
    
    \item[$\bullet$] \textbf{Valor social} (\textit{social value}): utilidad derivada de la capacidad del producto o servicio para mejorar la autoimagen, el estatus o el reconocimiento social  percibido por el cliente.
    
    \item[$\bullet$] \textbf{Valor de calidad/rendimiento} (\textit{quality/performance value}): utilidad derivada de la calidad percibida, la fiabilidad y el rendimiento esperado del producto o servicio servicio en relación con las necesidades y expectativas del usuario.
    
    \item[$\bullet$] \textbf{Valor del precio} (\textit{price/value for money}): utilidad derivada de  la percepción de que los beneficios obtenidos justifican el coste económico asumido, es decir, la sensación de obtener un retorno adecuado  por el dinero invertido.
\end{itemize}

La naturaleza multidimensional de PERVAL resulta especialmente adecuada para el análisis de procesos de negocio orientados al cliente. En muchos contextos organizacionales, los procesos no solo buscan alcanzar objetivos operativos de eficiencia y productividad, sino también generar experiencias que aporten valor a los usuarios finales. Sin embargo, el valor diseñado por una organización no siempre coincide con el valor efectivamente percibido durante la interacción real con el proceso.

En el contexto de este TFM, la clasificación PERVAL se emplea como marco conceptual  para analizar las actividades de un proceso de negocio modelado en BPMN e identificar de qué manera dichas actividades contribuyen a la generación de las distintas dimensiones  de valor percibido desde la perspectiva del cliente. Este análisis permite detectar posibles desalineamientos entre el valor  que la organización pretende ofrecer y el valor que efectivamente percibe el usuario durante su interacción con el proceso.

La incorporación de PERVAL en el presente trabajo adquiere una relevancia particular, ya que permite complementar la generación automática de modelos BPMN mediante LLMs con una perspectiva explícitamente orientada al cliente. De esta manera, el análisis no se limita únicamente a la representación estructural del proceso, sino que también considera las implicaciones que las diferentes actividades y decisiones del proceso pueden tener sobre la experiencia de valor del usuario final.

%%%%%%%%%%%%%%%%%%%%%%%%%%%%%%%%%%%%%%%%%%%%%%%%%%%%%%%%%%%
%% Final del capítulo de fundamentos
%%%%%%%%%%%%%%%%%%%%%%%%%%%%%%%%%%%%%%%%%%%%%%%%%%%%%%%%%%%
